\documentclass[aspectratio=169,11pt]{beamer}

% ============================================================
% THEME AND PACKAGES
% ============================================================

\usetheme{Singapore}
\usecolortheme{default}

\usepackage{graphicx}
\usepackage{booktabs}
\usepackage{tabularx}
\usepackage{array}
\usepackage{colortbl}
\usepackage{xcolor}
\usepackage{amsmath}
\usepackage{siunitx}

\graphicspath{{figures/}}

% ============================================================
% DRIP-FEED SWITCH
% ============================================================
%
% \dripfeedtrue  = reveal points gradually
% \dripfeedfalse = show complete slides
%
% ============================================================

\newif\ifdripfeed

\dripfeedtrue
% \dripfeedfalse

\ifdripfeed
  % Future items remain visible but faint.
  \setbeamercovered{transparent=18}

  \newcommand{\stepitem}{\item<+->}
  \newcommand{\reveal}[2]{\uncover<#1->{#2}}
\else
  \setbeamercovered{invisible}

  \newcommand{\stepitem}{\item}
  \newcommand{\reveal}[2]{#2}
\fi

% ============================================================
% GENERAL APPEARANCE
% ============================================================

\setbeamertemplate{navigation symbols}{}
\setbeamertemplate{caption}[numbered]

% IMPORTANT:
% Do not remove the Singapore headline.
% This preserves the soft shaded strip at the top.

% Slide title:
% - slightly lower than the previous version
% - minimal space beneath it
\setbeamertemplate{frametitle}{
  \nointerlineskip
  \vspace*{0.22cm}
  \hspace*{0.12cm}
  \insertframetitle\par
  \vspace*{-0.14cm}
}

% ============================================================
% COLOURS
% ============================================================

\definecolor{projectblue}{RGB}{28,76,128}
\definecolor{lightblue}{RGB}{232,240,248}
\definecolor{verylightblue}{RGB}{246,249,252}
\definecolor{tableStripe}{RGB}{239,245,250}
\definecolor{winnerRow}{RGB}{218,237,224}

% Pipeline stage colours.
\definecolor{stageone}{RGB}{224,237,249}
\definecolor{stagetwo}{RGB}{224,241,228}
\definecolor{stagethree}{RGB}{250,239,214}

\setbeamercolor{structure}{fg=projectblue}
\setbeamercolor{frametitle}{fg=projectblue}

% ============================================================
% GLOBAL FONT SIZE
% OPTION 2 — MODERATELY COMPACT
% ============================================================

\setbeamerfont{frametitle}{
  size=\fontsize{24}{25},
  series=\bfseries
}

\setbeamerfont{itemize/enumerate body}{
  size=\fontsize{17}{21}
}

\setbeamerfont{itemize/enumerate subbody}{
  size=\fontsize{14.5}{18.5}
}

% Title-page sizes.
\setbeamerfont{title}{
  size=\fontsize{31}{37},
  series=\bfseries
}

\setbeamerfont{author}{
  size=\fontsize{20}{25}
}

\setbeamerfont{institute}{
  size=\fontsize{17}{22}
}

\setbeamerfont{date}{
  size=\fontsize{16}{20}
}

\setlength{\leftmargini}{1.3em}

\newcommand{\bigbulletlist}{
  \setlength{\itemsep}{0.62em}
  \setlength{\parskip}{0pt}
}

% ============================================================
% TABLE SETTINGS
% ============================================================

\sisetup{
  table-format=2.2,
  table-number-alignment=center,
  detect-weight=true,
  detect-inline-weight=math
}

\renewcommand{\arraystretch}{1.28}
\setlength{\tabcolsep}{6pt}

% ============================================================
% TITLE INFORMATION
% ============================================================

\title[Restricted Dog Classification]
{Restricted Dog Breed Classification\\Using CNN Transfer Learning}

\author{Ronan Downes}

\institute{
  MSc in Data Analytics\\
  CCT College Dublin\\[0.25em]
  Supervisor: Sam Weiss
}

\date{April 2025}

% ============================================================
% DOCUMENT
% ============================================================

\begin{document}

% ============================================================
% TITLE SLIDE
% ============================================================

\begin{frame}

\centering

% Moves the title down slightly.
\vspace*{0.55cm}

{\usebeamerfont{title}
 \color{projectblue}
 {\Huge Restricted Dog Breed Classification\\[0.28em]
 Using CNN Transfer Learning
}}

\vspace{1.05em}
 {\LARGE
{\usebeamerfont{author}
 Ronan Downes
}

\vspace{0.78em}

{\usebeamerfont{institute}
 MSc in Data Analytics\\
 CCT College Dublin\\[0.20em]
 Supervisor: Sam Weiss
}}

\vspace{0.78em}

{\usebeamerfont{date}
 April 2025
}

\end{frame}

% ============================================================
% RESEARCH PROBLEM
% ============================================================

\begin{frame}{Research problem}

\vspace{-0.3em}

\begin{itemize}
  \bigbulletlist

  \stepitem Benchmark six pretrained CNNs
  \stepitem Binary dog-image classification
  \stepitem Irish restricted-breed context
  \stepitem Stanford Dogs dataset
  \stepitem Controlled model comparison
\end{itemize}

\end{frame}

% SAY ORALLY:
% One common transfer-learning pipeline was used to benchmark
% six ImageNet-pretrained CNN architectures.
%
% The classes were informed by the Irish restricted-dog
% regulations, but this was a controlled image-classification
% experiment rather than legal breed identification.

% ============================================================
% RESEARCH QUESTIONS
% ============================================================

\begin{frame}{Research questions}

\vspace{-0.3em}

\begin{enumerate}
  \bigbulletlist

  \stepitem How do six frozen CNNs compare?
  \stepitem Does fine-tuning improve performance?
  \stepitem What does Grad-CAM reveal?
  \stepitem What limits real-world use?
\end{enumerate}

\end{frame}

% ============================================================
% RESTRICTED-BREED COVERAGE
% ============================================================

\begin{frame}{Stanford Dogs Restricted-breeds}

\vspace{-0.35em}

\centering

{\fontsize{15.5}{19}\selectfont

\begin{tabularx}{0.92\textwidth}{
  >{\raggedright\arraybackslash}X
  >{\raggedright\arraybackslash}X
}

\textbf{\color{projectblue}Represented}
&
\textbf{\color{projectblue}Not represented}
\\[-0.15em]

\cmidrule(r){1-1}
\cmidrule(l){2-2}

Rottweiler
&
{American Pit Bull Terrier}
\\

\rowcolor{tableStripe}
German Shepherd
&
English Bull Terrier
\\

Doberman Pinscher
&
Japanese Akita
\\

\rowcolor{tableStripe}
Bull Mastiff
&
Japanese Tosa
\\

Rhodesian Ridgeback
&
Bandog
\\

\rowcolor{tableStripe}
Staffordshire Bull Terrier
&
\\

\end{tabularx}
}

\end{frame}

% ============================================================
% DATASET CONSTRUCTION
% ============================================================

\begin{frame}{Dataset construction}

\vspace{-0.25em}

\begin{columns}[c,onlytextwidth]

\column{0.38\textwidth}

\centering

{\fontsize{36}{40}\selectfont\bfseries 6}

\vspace{0.08em}

{\large restricted breeds}

\column{0.12\textwidth}

\centering
{\Large\bfseries vs.}

\column{0.38\textwidth}

\centering

{\fontsize{36}{40}\selectfont\bfseries 24}

\vspace{0.08em}

{\large unrestricted breeds}

\end{columns}

\vspace{0.6em}

\centering

{\normalsize
24 breeds reproducibly selected\\
from the remaining 114 Stanford categories
}

\vspace{0.55em}

\begin{center}
\begin{minipage}{0.68\textwidth}

\begin{itemize}
  \setlength{\itemsep}{0.42em}

  \stepitem Broader negative class
  \stepitem Manageable computational scope
  \stepitem Fixed data splits
\end{itemize}

\end{minipage}
\end{center}

\end{frame}

% SAY ORALLY:
% Stanford Dogs contains 120 breed categories.
% Six mapped to represented restricted breeds.
% The unrestricted class was reproducibly sampled from the
% remaining 114 categories.
%
% This provided a broader negative class while keeping the
% benchmark computationally manageable.

% ============================================================
% EXPERIMENTAL PIPELINE
% ============================================================

\begin{frame}{Experimental pipeline}

\vspace{-0.25em}

\centering

% ------------------------------------------------------------
% STAGE 1 — FROZEN BENCHMARK
% First four cells
% ------------------------------------------------------------

\begin{columns}[T,onlytextwidth]

\column{0.235\textwidth}
\reveal{1}{
\centering
\fcolorbox{projectblue}{stageone}{
  \parbox[c][1.55cm][c]{0.88\linewidth}{
    \centering
    \normalsize\bfseries
    Dataset\\construction
  }
}
}

\column{0.235\textwidth}
\reveal{2}{
\centering
\fcolorbox{projectblue}{stageone}{
  \parbox[c][1.55cm][c]{0.88\linewidth}{
    \centering
    \normalsize\bfseries
    Resize and\\normalise
  }
}
}

\column{0.235\textwidth}
\reveal{3}{
\centering
\fcolorbox{projectblue}{stageone}{
  \parbox[c][1.55cm][c]{0.88\linewidth}{
    \centering
    \normalsize\bfseries
    Training-data\\augmentation
  }
}
}

\column{0.235\textwidth}
\reveal{4}{
\centering
\fcolorbox{projectblue}{stageone}{
  \parbox[c][1.55cm][c]{0.88\linewidth}{
    \centering
    \normalsize\bfseries
    Frozen CNN\\benchmark
  }
}
}

\end{columns}

\vspace{0.75em}

% ------------------------------------------------------------
% STAGE 2 AND STAGE 3
% ------------------------------------------------------------

\begin{columns}[T,onlytextwidth]

\column{0.22\textwidth}

\column{0.26\textwidth}
\reveal{5}{
\centering
\fcolorbox{projectblue}{stagetwo}{
  \parbox[c][1.55cm][c]{0.88\linewidth}{
    \centering
    \normalsize\bfseries
    Fine-tune\\selected CNN
  }
}
}

\column{0.26\textwidth}
\reveal{6}{
\centering
\fcolorbox{projectblue}{stagethree}{
  \parbox[c][1.55cm][c]{0.88\linewidth}{
    \centering
    \normalsize\bfseries
    Grad-CAM\\explainability
  }
}
}

\column{0.22\textwidth}

\end{columns}

\vspace{0.5em}

\reveal{7}{
\centering

{\small
\colorbox{stageone}{\strut Stage 1: frozen benchmark}
\hspace{0.7em}
\colorbox{stagetwo}{\strut Stage 2: fine-tuning}
\hspace{0.7em}
\colorbox{stagethree}{\strut Stage 3: explainability}
}
}

\end{frame}

% SAY ORALLY:
% Images were resized to the input resolution required by each
% architecture: either 224 x 224 or 299 x 299.
%
% Each architecture's own preprocess_input function was used.
%
% Training augmentation included horizontal flipping,
% small rotations and small zoom transformations.

% ============================================================
% SHORT ARCHITECTURE TIMELINE
% ============================================================

\begin{frame}{ILSVRC Winners Record}

\vspace{-0.37em}

\centering

{\fontsize{14.5}{17.5}\selectfont

\begin{tabularx}{0.92\textwidth}{
  >{\centering\arraybackslash}p{0.12\textwidth}
  >{\raggedright\arraybackslash}p{0.30\textwidth}
  >{\raggedright\arraybackslash}X
}

\textbf{\color{projectblue}Year}
&
\textbf{\color{projectblue}Architecture}
&
\textbf{\color{projectblue}Key idea}
\\[-0.15em]

\cmidrule(r){1-1}
\cmidrule(lr){2-2}
\cmidrule(l){3-3}

2010
&
SIFT/LBP + SVM
&
Hand-crafted features 
\\

\rowcolor{tableStripe}
2011
&
Fisher vectors
&
Best pre-CNN 
\\

2012
&
AlexNet
&
Deep-CNN breakthrough
\\

2014
&
VGG16
&
Uniform deep CNN
\\

\rowcolor{tableStripe}
2015
&
ResNet50
&
Residual connections
\\

2015
&
InceptionV3
&
Multi-scale processing
\\

\rowcolor{tableStripe}
2016
&
Xception
&
Separable convolutions
\\

2016
&
InceptionResNetV2
&
Inception plus residual
\\

\rowcolor{tableStripe}
2017
&
NASNetMobile
&
Neural architecture search
\end{tabularx}
}




\end{frame}

% ============================================================
% SIX BENCHMARKED CNNs
% ============================================================

\begin{frame}{Six benchmarked CNN models}

\vspace{-0.25em}

\begin{columns}[T,onlytextwidth]

\column{0.48\textwidth}

\begin{itemize}
  \setlength{\itemsep}{0.7em}

  \stepitem \textbf{VGG16}\\
        {\small Uniform deep baseline}

  \stepitem \textbf{ResNet50}\\
        {\small Residual connections}

  \stepitem \textbf{InceptionV3}\\
        {\small Multi-scale processing}
\end{itemize}

\column{0.48\textwidth}

\begin{itemize}
  \setlength{\itemsep}{0.7em}

  \stepitem \textbf{Xception}\\
        {\small Separable convolutions}

  \stepitem \textbf{InceptionResNetV2}\\
        {\small Multi-scale residual model}

  \stepitem \textbf{NASNetMobile}\\
        {\small Architecture-search model}
\end{itemize}

\end{columns}

\end{frame}

% ============================================================
% FROZEN BENCHMARK METHOD
% ============================================================

\begin{frame}{Frozen Benchmark Approach}

\vspace{-0.25em}

\begin{columns}[T,onlytextwidth]

\column{0.48\textwidth}

\begin{itemize}
  \setlength{\itemsep}{0.63em}

  \stepitem ImageNet weights
  \stepitem Frozen feature extractors
  \stepitem Common classification head
\end{itemize}

\column{0.48\textwidth}

\begin{itemize}
  \setlength{\itemsep}{0.63em}

  \stepitem Class weighting
  \stepitem Early stopping
  \stepitem Best validation checkpoint
\end{itemize}

\end{columns}

\reveal{7}{
\vspace{0.4em}

\centering
{\large\bfseries Multiple metrics—not accuracy alone}
}

\end{frame}


% ============================================================
% BENCHMARK SCOPE
% ============================================================


% ============================================================
% FROZEN RESULTS — NATIVE LATEX TABLE
% ============================================================

\begin{frame}{Frozen Benchmark Results Table}

\vspace{-0.4em}

\centering

{\fontsize{16.2}{19.5}\selectfont

\rowcolors{2}{white}{tableStripe}

\begin{tabularx}{0.97\textwidth}{
  >{\raggedright\arraybackslash}p{0.37\textwidth}
  S[table-column-width=0.12\textwidth]
  S[table-column-width=0.12\textwidth]
  S[table-column-width=0.12\textwidth]
  S[table-column-width=0.10\textwidth]
}
\textbf{\color{projectblue}CNN}
&
{\textbf{\color{projectblue}Acc.}}
&
{\textbf{\color{projectblue}Prec.}}
&
{\textbf{\color{projectblue}Recall}}
&
{\textbf{\color{projectblue}F1}}
\\[-0.12em]

\cmidrule(r){1-1}
\cmidrule(lr){2-2}
\cmidrule(lr){3-3}
\cmidrule(lr){4-4}
\cmidrule(l){5-5}

VGG16
& 87.4
& 60.3
& 92.1
& 72.9
\\

ResNet50
& 91.7
& 70.9
& 93.7
& 80.7
\\

InceptionV3
& 96.1
& 84.4
& 96.8
& 90.2
\\

Xception
& 96.0
& 83.7
& 97.4
& 90.0
\\

\rowcolor{winnerRow}
\textbf{InceptionResNetV2}
& \bfseries 96.8
& \bfseries 86.9
& \bfseries 97.4
& \bfseries 91.8
\\

NASNetMobile
& 94.8
& 79.9
& 96.3
& 87.4
\\

\end{tabularx}
}

\vspace{0.22em}

{\scriptsize Percentages shown; bold and shading are selected model.}

\end{frame}

% ============================================================
% FROZEN BAR CHART
% ============================================================

\begin{frame}{Frozen Benchmark Results Chart}

\vspace{-0.35em}

\centering

\includegraphics[
  width=0.99\textwidth,
  height=0.81\textheight,
  keepaspectratio
]{model_zoo_four_metric_comparison_presentation.pdf}

\end{frame}

% ============================================================
% MODEL SELECTION
% ============================================================

% ============================================================
% FINE-TUNING STRATEGY
% ============================================================

\begin{frame}{Fine-tuning strategy}

\vspace{-0.3em}

\centering

{\normalsize\bfseries
Task adaptation without destabilising pretrained features
}

\vspace{0.55em}

\begin{columns}[T,onlytextwidth]

\column{0.47\textwidth}

\begin{itemize}
  \setlength{\itemsep}{0.62em}

  \stepitem Last 20 layers unfrozen
  \stepitem Learning rate: \(10^{-5}\)
  \stepitem Maximum 15 epochs
\end{itemize}

\column{0.47\textwidth}

\begin{itemize}
  \setlength{\itemsep}{0.62em}

  \stepitem Class weighting retained
  \stepitem Early stopping
  \stepitem Best validation checkpoint
\end{itemize}

\end{columns}

\end{frame}

% ============================================================
% EFFECT OF FINE-TUNING
% ============================================================

\begin{frame}{Fine-tuning results}

\vspace{-0.3em}

\begin{columns}[T,onlytextwidth]

\column{0.4\textwidth}

\centering

{\small

\renewcommand{\arraystretch}{1.22}

\begin{tabular}{lcc}
\toprule
Metric & Frozen & Fine-tuned\\
\midrule
Accuracy  & 96.79\% & \textbf{98.25\%}\\
Precision & 86.85\% & \textbf{93.88\%}\\
Recall    & \textbf{97.37\%} & 96.84\%\\
F1        & 91.81\% & \textbf{95.34\%}\\
MCC       & 0.9005 & \textbf{0.9428}\\
\bottomrule
\end{tabular}
}

\vspace{0.5em}

{\normalsize\bfseries
Fewer false positives\\
Slightly more false negatives
}

\column{0.6\textwidth}

\centering

\includegraphics[
  width=\linewidth,
  height=0.84\textheight,
  keepaspectratio
]{FineTuned_InceptionResNetV2_confusion_matrix.png}

\end{columns}

\end{frame}

% ============================================================
% LEARNING BEHAVIOUR
% ============================================================

\begin{frame}{Learning behaviour}

\vspace{-0.35em}

\centering

\includegraphics[
  width=0.94\textwidth,
  height=0.77\textheight,
  keepaspectratio
]{FineTuned_InceptionResNetV2_learning_curves.png}

{\small\bfseries Best validation checkpoint retained}

\end{frame}

% ============================================================
% GRAD-CAM — ALL FOUR SHOWN
% ============================================================

\begin{frame}{Grad-CAM explainability}

\vspace{-0.35em}

\centering

% Top row
\begin{columns}[T,onlytextwidth]

\column{0.48\textwidth}

\centering

\includegraphics[
  width=0.82\linewidth,
  height=0.29\textheight,
  keepaspectratio
]{TP_rottweiler_n02106550_4962_gradcam.png}

{\small True positive}

\column{0.48\textwidth}

\centering

\includegraphics[
  width=0.82\linewidth,
  height=0.29\textheight,
  keepaspectratio
]{TN_kelpie_n02105412_7810_gradcam.png}

{\small True negative}

\end{columns}

\vspace{0.28em}

% Bottom row
\begin{columns}[T,onlytextwidth]

\column{0.48\textwidth}

\centering

\includegraphics[
  width=0.82\linewidth,
  height=0.29\textheight,
  keepaspectratio
]{FP_weimaraner_n02092339_3028_gradcam.png}

{\small False positive}

\column{0.48\textwidth}

\centering

\includegraphics[
  width=0.82\linewidth,
  height=0.29\textheight,
  keepaspectratio
]{FN_staffordshire_bullterrier_n02093256_5988_gradcam.png}

{\small False negative}

\end{columns}

\vspace{0.12em}

{\scriptsize Diagnostic evidence—not proof of causal reasoning.}

\end{frame}



% ============================================================
% REPRODUCIBILITY
% ============================================================
% ============================================================
% PROJECT NOTEBOOKS AND EXECUTION ENVIRONMENT
% ============================================================
% ============================================================
% COLAB NOTEBOOKS
% ============================================================


\begin{frame}{Project Contribution}

\vspace{-0.3em}

\begin{columns}[T,onlytextwidth]

\column{0.47\textwidth}

\begin{itemize}
  \setlength{\itemsep}{0.65em}

  \stepitem Reproducible pipeline
  \stepitem Six-CNN benchmarks
  \stepitem Multi-metric model selection
\end{itemize}

\column{0.47\textwidth}

\begin{itemize}
  \setlength{\itemsep}{0.65em}

  \stepitem Fine-tuning analysis
  \stepitem Error trade-off analysis
  \stepitem Grad-CAM explainability
\end{itemize}

\end{columns}

\end{frame}

\begin{frame}{IPYNB Notebook Execution}

\vspace{-0.25em}

\begin{enumerate}
\setlength{\itemsep}{0.62em}

\stepitem \texttt{01\_Data\_Preparation.ipynb}
\hfill {\small\color{projectblue}\textbf{CPU}}

\stepitem \texttt{02\_Benchmarking.ipynb}
\hfill {\small\color{projectblue}\textbf{GPU required}}

\stepitem \texttt{03\_Fine\_Tuning.ipynb}
\hfill {\small\color{projectblue}\textbf{GPU required}}

\stepitem \texttt{04\_Gradcam.ipynb}
\hfill {\small\color{projectblue}\textbf{CPU}}

\end{enumerate}

\vspace{0.45em}

\centering
{\small
Local hardware did not provide a suitable GPU, so Google Colab was used to
accelerate benchmarking and fine-tuning with a \textbf{GPU hardware accelerator}
(NVIDIA T4 when available). Free-tier compute limits occasionally restricted
GPU access, so some runs were completed during off-peak periods.
}

\end{frame}
% ============================================================
% CONCLUSIONS
% ============================================================
% ============================================================
% LIMITATIONS — VIVA PROMPT
% ============================================================
% \begin{frame}[fragile]{GPU Runtime Configuration}

% \vspace{-0.3em}

% \begin{block}{Google Colab runtime}
% \small
% Benchmarking and fine-tuning were executed in Google Colab using a
% \textbf{GPU hardware accelerator}. When available, Colab assigned an
% \textbf{NVIDIA T4 GPU}.
% \end{block}

% \vspace{0.35em}

% \begin{block}{TensorFlow GPU check}
% \begin{verbatim}
% import tensorflow as tf

% gpus = tf.config.list_physical_devices("GPU")

% if not gpus:
%     raise RuntimeError("No GPU detected")
% \end{verbatim}
% \end{block}

% \vspace{0.3em}

% \centering
% {\small
% The notebook did not request a specific GPU model;
% it verified that a GPU had been assigned to the Colab runtime.
% }

% \end{frame}
\begin{frame}

\centering
\vfill

{\fontsize{40}{46}\selectfont
 \color{projectblue}
 \bfseries
 Limitations?
}

\vfill

\end{frame}

% TALK ABOUT:
% - only 6 of 11 restricted categories represented
% - no mixed ancestry
% - image-only evidence
% - breed-level variation hidden by aggregate metrics
% - dataset bias and distribution shift

% ============================================================
% ETHICAL AND LEGAL BOUNDARY — VIVA PROMPT
% ============================================================

\begin{frame}

\centering
\vfill

{\fontsize{36}{43}\selectfont
 \color{projectblue}
 \bfseries
Deployment?
}

\vfill

\end{frame}



\begin{frame}

\centering
\vfill

{\fontsize{36}{43}\selectfont
 \color{projectblue}
 \bfseries
Further Questions?
}

\vfill

\end{frame}
% TALK ABOUT:
% - statistical class is not ancestry
% - class label is not dangerousness
% - photographs cannot establish legal status
% - false positives and false negatives have different consequences
% - research prototype, not an enforcement system
\begin{frame}{Conclusions}

\vspace{-0.3em}

\begin{columns}[T,onlytextwidth]

\column{0.49\textwidth}

\begin{enumerate}
  \setlength{\itemsep}{0.55em}

  \stepitem Transfer learning performed strongly
  \stepitem InceptionResNetV2 ranked first
  \stepitem Fine-tuning improved overall balance
\end{enumerate}

\column{0.45\textwidth}

\begin{enumerate}
  \setcounter{enumi}{3}
  \setlength{\itemsep}{0.55em}

  \reveal{4}{\item Fine-tuning changed the error trade-off}
  \reveal{5}{\item Current limitations do not justify real-world deployment}
\end{enumerate}

\end{columns}



\end{frame}
% ============================================================
% QUESTIONS
% ============================================================



\end{document}
